\fichatitulo{15 · APLICAÇÃO}{Artigo · 2025}{LegisSearch}
{\small Colombo et al. \textit{LegisSearch: navigating legislation with graphs and large language models}. Artigo · 2025. \href{https://link.springer.com/content/pdf/10.1007/s10506-025-09482-6.pdf}{Fonte consultada}.\par}

\campo{Problema e objetivo}
Como melhorar a localização de legislação italiana relevante?

\campo{Principal contribuição · síntese interpretativa}
Integra grafo de propriedades, metadados legislativos, representações vetoriais e etapas com LLM para apoiar busca e navegação.

\campo{Método / natureza da evidência}
Comparação da recuperação com e sem expansão e filtragem por LLM, usando conjuntos temáticos de atos relevantes.

\campo{Resultado ou argumento principal}
A média de recall@5 passa de 0,25 para 0,35 na tabela 6; os ganhos variam entre temas. Trata-se de recuperação, não de avaliação de respostas geradas.

\campo{Excertos-chave · seleção literal}
\begin{itemize}
\excerto{1}{we do not have a ranking of such laws}{Discussão das limitações.}
\excerto{2}{the impact of LLM-based query expansion varies across thematic areas}{Discussão da tabela 6.}
\end{itemize}

\campo{Limitações e análise crítica}
Autores: o conjunto de referência não ordena as leis relevantes. Análise da ficha: consultas oriundas de listas temáticas podem diferir das necessidades de usuários reais.

\campo{Aproveitamento na dissertação · proposta}
Trabalhos relacionados: comparar recuperação e filtros de metadados, com análise por tema.

\registro{Fonte consultada nesta preparação: Resumo, introdução, tabela 6 e discussão das limitações. Chave: \texttt{colomboEtAl2025legisSearch}.}
